\section{模型假设}

本文的建模建立在下列假设之上。
每一条假设的影响都在第 6 节中量化。

\begin{enumerate}[label=\arabic*., leftmargin=2.6em, itemsep=4pt]
  \item \textbf{局部热力学平衡}：药材内部各点满足局部平衡，烘房环境的温度与水分浓度可以作为
        边界驱动量；否则第三类边界条件无法建立。
  \item \textbf{一维径向轴对称}：药材为长径比大于 $12$ 的圆柱，热风沿轴向与周向均匀作用，
        故只保留径向坐标，忽略两端面的轴向输运。该简化在 $30$~min 的预热窗口内影响不足长度的 $2\%$，
        但在数日的干燥时程内轴向渗透可达数厘米每端，其影响在检验部分量化。
  \item \textbf{质量方程取标准的 Fick 形式}：按题目明文的写法定式，密度、比热容与导热系数只进入能量方程。
        若改用与湿基密度严格自洽的写法，等效扩散系数会相差一个随浓度变化的因子，其对结果的影响在检验部分量化。
  \item \textbf{物性与经验式取题目给定值}：密度、比热容、导热系数、扩散系数与对流传热传质系数
        均按题目附录取值，不作拟合。
  \item \textbf{恒温干燥段的环境为定值}：附件 1 只给到 $14400$~s，其后按已确认的口径
        维持恒定的环境温度与水分浓度；该设定由附件 1 稳定段的实测均值核实。
  \item \textbf{忽略水分汽化潜热}：题目未给潜热数据，能量方程中不含相变项，
        该简化会使升温阶段的预测偏快。
  \item \textbf{变物性系数取上一时间层}：使每个时间步仍是线性方程组，代价是时间格式降为一阶，
        其精度由时间加密实验核验。
  \item \textbf{收缩各向同性且严格按附件 2}：问题四的半径随时间的变化完全由附件 2 决定，
        收缩只改变几何尺寸，不改变干基含水率的定义。
  \item \textbf{问题四的报告坐标取当前物理距离}：按题目已确认的口径，
        每个时刻按当时半径换算物理坐标，超出当前半径的位置不给出数值。
  \item \textbf{问题四的表观对流系数由附件 2 数值微分得到}：其噪声放大风险由\textbf{解析不变量检验}控制。
\end{enumerate}
